\chapter{LANDASAN KEPUSTAKAAN}

Bab ini membahas penelitian terdahulu yang berkaitan dengan prediksi harga emas serta teori yang digunakan dalam penelitian. Kajian pustaka menunjukkan posisi penelitian berdasarkan objek, target, fitur, model, dan evaluasi. Landasan teori menjelaskan konsep-konsep yang mendasari pembentukan tiga konfigurasi fitur dan pengukuran perubahan kinerjanya.

\section{Kajian Pustaka}

Kajian penelitian mutakhir menggunakan artikel jurnal final tahun 2021--2026, sedangkan rujukan yang lebih lama digunakan sebagai dasar teori. Studi yang dibahas mencakup klasifikasi arah, hubungan makroekonomi, penggabungan fitur, model prediksi, dan validasi temporal.

\subsection{Klasifikasi Arah dan Representasi Teknikal}

\textcite{sadorsky2021predicting} mengklasifikasikan arah harian ETF emas GLD dan ETF perak SLV untuk horizon 1--20 hari menggunakan 13 indikator teknikal. Penelitian tersebut membandingkan model logit, bagging, boosting, dan random forest. Model berbasis pohon memberikan hasil yang lebih baik daripada model linear pada sebagian besar konfigurasi pengujian. Studi ini menjadi dasar penggunaan model dengan karakter linear dan berbasis pohon, sedangkan penelitian yang diusulkan menggunakan harga spot XAUUSD H1 dan target tiga kelas.

\textcite{basher2022forecasting} memprediksi arah harian emas dan Bitcoin dengan fitur teknikal dan makroekonomi, termasuk suku bunga, inflasi, dan volatilitas pasar. Model logit, bagging, dan random forest dievaluasi menggunakan validasi silang deret waktu serta analisis kepentingan variabel. Penelitian tersebut menunjukkan bahwa fitur makroekonomi dapat disertakan dalam klasifikasi arah dan dievaluasi secara temporal. Penelitian ini melanjutkan gagasan tersebut pada klasifikasi BUY, HOLD, dan SELL dalam timeframe satu jam.

\textcite{mamoudan2023hybrid} mengklasifikasikan sinyal teknikal sebagai \emph{FAKE} atau \emph{REAL}. Arsitekturnya menggabungkan CNN--BiGRU, pemilihan fitur dengan Moth-Flame Optimization, dan penyetelan dengan Firefly Algorithm. Kandidat masukannya mencakup MACD, \emph{moving average}, dan Ichimoku. Studi tersebut memperlihatkan pemanfaatan transformasi teknikal, tetapi targetnya berupa validitas sinyal dan bukan arah harga berikutnya. Dalam penelitian ini, indikator teknikal digunakan untuk merepresentasikan momentum dan tren dari riwayat harga agar dapat dipelajari oleh model linear maupun nonlinear.

\subsection{Hubungan Makroekonomi dan Informasi Intrahari}

Hubungan antara pasar emas dan kondisi ekonomi telah dikaji melalui beberapa pendekatan ekonometrika. \textcite{dinh2022economic} menggunakan GARCH-MIDAS dan DCC-MIDAS untuk menganalisis volatilitas dan korelasi logam mulia dalam kaitannya dengan variabel makroekonomi. \textcite{wang2023asymmetric} menggunakan regresi kuantil untuk menilai hubungan asimetris harga minyak dan dolar Amerika Serikat dengan imbal hasil logam mulia, disertai federal funds rate dan VIX sebagai variabel kontrol. Hasil kedua studi menunjukkan bahwa hubungan tersebut dapat berbeda menurut kondisi pasar dan bagian distribusi yang diamati. Oleh karena itu, manfaat fitur makroekonomi untuk prediksi perlu diukur pada tugas dan periode pengujian yang ditetapkan.

Pada resolusi intrahari, \textcite{sobti2021how} menganalisis data satu menit periode 2013--2018 dan 62 pengumuman makroekonomi untuk mempelajari pembentukan harga di pasar emas. Temuan tersebut menunjukkan bahwa informasi pengumuman ekonomi relevan sepanjang hari perdagangan. \textcite{yousaf2022information} juga menggunakan data per jam untuk mengkaji limpahan informasi dan efektivitas lindung nilai pasangan minyak--emas serta minyak--Bitcoin selama COVID-19. Kedua studi ini mendukung penggunaan resolusi intrahari dan pengujian informasi eksternal, meskipun kontribusinya terhadap klasifikasi XAUUSD H1 tetap harus diukur secara \emph{out-of-sample}.

\subsection{Penggabungan Fitur, Model, dan Validasi Temporal}

\textcite{taneva2025framework} memadukan data finansial, makroekonomi, dan sentimen dengan metode ekonometrika, pembelajaran mesin, dan pembelajaran mendalam untuk meramalkan harga emas harian periode 2014--2024. Penelitian tersebut menggunakan target harga kontinu dan metrik regresi. Selain perbedaan target, penggunaan fitur berfrekuensi rendah pada keputusan H1 memerlukan pemeriksaan waktu ketersediaan data agar informasi masa depan tidak masuk ke fitur.

\textcite{velez2026predictive} meramalkan harga emas kontinu bulanan periode Januari 2020--Juni 2026 melalui 32 kombinasi model dan fitur dengan validasi temporal. Studi tersebut memberikan preseden langsung bagi perbandingan konfigurasi fitur. Penelitian ini menggunakan rancangan yang berbeda, yaitu tiga konfigurasi harga--teknikal--makro yang ditambahkan secara bertahap untuk klasifikasi tiga arah H1 pada bar dan periode uji yang sama.

\textcite{hisyam2025implementasi} menggunakan XGBoost, optimasi Bayesian, dan \emph{sliding walk-forward validation} 14 hari untuk meramalkan harga penutupan dan \emph{buyback} emas Antam harian dalam rupiah per gram. RSI, EMA, dan MACD dianalisis melalui kepentingan fitur. Sementara itu, \textcite{arif2026multivariate} menggunakan LSTM untuk meramalkan harga penutupan XAUUSD H1 tahun 2024 berdasarkan data OHLC dan lima pemisahan \emph{rolling-origin}. Kedua studi mendukung penggunaan XGBoost, indikator teknikal, objek XAUUSD, dan evaluasi temporal pada konteks masing-masing. Penelitian ini berfokus pada perubahan skor antarkelompok fitur setelah model dilatih dan disetel ulang untuk klasifikasi tiga kelas.

\subsection{Perbandingan Studi Terdekat dan Posisi Penelitian}

Tabel \ref{tab:studi-terdekat} merangkum tujuh penelitian yang paling dekat berdasarkan objek, target, fitur, model, atau cara evaluasinya. Perbedaan yang disajikan memperjelas posisi penelitian ini terhadap tujuan setiap studi.

\begingroup
\setlength{\LTleft}{\fill}
\setlength{\LTright}{\fill}
\setlength{\tabcolsep}{3pt}
\begin{longtable}{>{\raggedright\arraybackslash}p{2.2cm}>{\raggedright\arraybackslash}p{2.8cm}>{\raggedright\arraybackslash}p{3.4cm}>{\raggedright\arraybackslash}p{4.2cm}}
    \caption{Perbandingan tujuh penelitian terdekat}
    \label{tab:studi-terdekat}\\
    \hline
    Studi & Target dan frekuensi & Metode dan fitur utama & Perbedaan dengan penelitian ini \\
    \hline
    \endfirsthead
    \caption[]{Perbandingan tujuh penelitian terdekat (lanjutan)}\\
    \hline
    Studi & Target dan frekuensi & Metode dan fitur utama & Perbedaan dengan penelitian ini \\
    \hline
    \endhead
    \hline
    \multicolumn{4}{r}{Bersambung ke halaman berikutnya}\\
    \endfoot
    \hline
    \endlastfoot
    Sadorsky (2021) & Arah biner ETF GLD/SLV harian; horizon 1--20 hari & Logit, bagging, boosting, random forest; 13 indikator teknikal & Harga spot XAUUSD H1; tiga kelas; evaluasi \emph{walk-forward}. \\
    Basher dan Sadorsky (2022) & Arah emas dan Bitcoin harian & Logit, bagging, random forest; fitur teknikal dan makro; validasi deret waktu & BUY/HOLD/SELL H1; tiga kelompok fitur bertahap. \\
    Mamoudan dkk. (2023) & Sinyal \emph{FAKE}/\emph{REAL}; frekuensi tidak dinyatakan & CNN--BiGRU dan metaheuristik; MACD, MA, Ichimoku & Arah harga satu bar berikutnya; dua model dengan karakter berbeda. \\
    Taneva-Angelova dkk. (2025) & Harga emas kontinu harian; 2014--2024 & Fusi finansial, makro, sentimen; ekonometrika, ML, DL & Klasifikasi H1; macro-F1 dan balanced accuracy. \\
    Velez Flores dkk. (2026) & Harga emas kontinu bulanan; Jan. 2020--Jun. 2026 & 32 kombinasi model--fitur; validasi temporal & Tiga kelas H1; penambahan kelompok pada bar uji yang sama. \\
    Hisyam dkk. (2025) & Harga Antam harian dalam rupiah per gram & XGBoost, optimasi Bayesian, \emph{sliding walk-forward}; RSI, EMA, MACD & XAUUSD H1; selisih skor kelompok fitur. \\
    Arif dkk. (2026) & Harga penutupan XAUUSD H1 tahun 2024 & LSTM multivariat, OHLC, lima pemisahan \emph{rolling-origin} & Klasifikasi tiga kelas; penambahan teknikal dan makro; dua model. \\
\end{longtable}
\endgroup

Dalam batas studi yang ditelaah, unsur klasifikasi arah, fitur teknikal dan makroekonomi, perbandingan konfigurasi fitur, XGBoost, XAUUSD H1, serta validasi temporal telah digunakan dalam penelitian terdahulu. Celah yang dikaji terletak pada penggabungan unsur-unsur tersebut untuk klasifikasi BUY, HOLD, dan SELL XAUUSD H1. Penelitian ini akan membandingkan tiga konfigurasi fitur kumulatif menggunakan model linear dan nonlinear pada observasi serta periode uji kronologis yang sama.

Kontribusi fitur akan diukur sebagai perubahan kinerja yang bersifat kondisional terhadap urutan penambahan kelompok, model, periode, dan penyetelan. Pada analisis utama, setiap konfigurasi akan disetel secara terpisah sehingga selisih skor turut mencerminkan respons penyetelan ulang. Analisis parameter tetap pada Bab 3 akan digunakan untuk memeriksa sensitivitas hasil. Dengan demikian, hasil penelitian akan dibaca sebagai temuan prediktif dan deskriptif dalam ruang lingkup pengujian.

\section{Landasan Teori}

\subsection{XAUUSD, Bar H1, dan Waktu Keputusan}

XAUUSD menyatakan harga satu troy ounce emas dalam dolar Amerika Serikat. Kenaikan XAUUSD menunjukkan penguatan emas relatif terhadap dolar Amerika Serikat, sedangkan penurunan menunjukkan keadaan sebaliknya. Data H1 merangkum satu jam perdagangan melalui harga pembukaan $O_t$, tertinggi $H_t$, terendah $L_t$, dan penutupan $C_t$. MetaTrader 5 juga dapat menyediakan \emph{tick volume}, yaitu jumlah tick atau pembaruan kuotasi yang tercatat, yang secara konseptual berbeda dari volume transaksi pada bursa terpusat. Kolom yang akan digunakan ditetapkan berdasarkan hasil pemeriksaan data pada Bab 3.

Pada akhir bar ke-$t$, riwayat sampai $C_t$ telah tersedia untuk membentuk fitur, sedangkan $C_{t+1}$ baru diketahui setelah bar berikutnya selesai. Pemisahan waktu fitur, keputusan, dan akhir horizon target mencegah kebocoran informasi masa depan. Imbal hasil logaritmik pada bar ke-$t$ dihitung menggunakan Persamaan \ref{eq:logreturn} untuk merepresentasikan perubahan harga secara relatif \parencite{tsay2010analysis}.
\begin{equation}
r_t=\ln\left(\frac{C_t}{C_{t-1}}\right)
\label{eq:logreturn}
\end{equation}
Pada persamaan tersebut, $r_t$ adalah imbal hasil logaritmik bar ke-$t$, $C_t$ adalah harga penutupan bar ke-$t$, dan $C_{t-1}$ adalah harga penutupan bar sebelumnya. Besaran relatif lain, seperti bentuk bar $\ln(C_t/O_t)$, rentang $\ln(H_t/L_t)$, dan volatilitas historis, digunakan untuk merangkum dinamika harga tanpa bergantung langsung pada tingkat harga. Seluruh jendela perhitungan berakhir paling lambat pada waktu keputusan.

\subsection{Target Tiga Kelas dan Ambang Historis}

Target penelitian memetakan imbal hasil satu bar berikutnya ke dalam kelas BUY, HOLD, atau SELL. Ketiga istilah yang lazim dalam konteks pasar tersebut digunakan sebagai label arah yang dibentuk dari harga, bukan label bawaan penyedia data maupun instruksi transaksi. Dalam pemodelan pembelajaran mesin finansial, penggunaan ambang batas dinamis berbasis volatilitas lokal bertujuan memisahkan fluktuasi tren yang bermakna dari derau acak pasar \parencite{lopezdeprado2018advances}. Dengan ambang volatilitas $\delta_t$, aturan pembentukan kelas ditunjukkan pada Persamaan \ref{eq:labeling}.
\begin{equation}
y_t=\begin{cases}\mathrm{BUY},&r_{t+1}>\delta_t\\\mathrm{HOLD},&-\delta_t\le r_{t+1}\le\delta_t\\\mathrm{SELL},&r_{t+1}<-\delta_t\end{cases}
\label{eq:labeling}
\end{equation}
Pada persamaan tersebut, $y_t$ adalah label untuk keputusan pada bar ke-$t$, $r_{t+1}$ adalah imbal hasil bar berikutnya, dan $\delta_t$ adalah ambang yang dihitung hanya dari data historis sampai bar ke-$t$. Kelas HOLD merepresentasikan perubahan yang kecil relatif terhadap volatilitas historis. Horizon dan rumus operasional ambang dijelaskan pada Bab 3.

Seluruh konfigurasi akan menggunakan label dan bar layak yang sama. Masker observasi bersama diperlukan agar selisih skor mencerminkan perubahan fitur, bukan perubahan sampel. Ketentuan yang sama berlaku pada kedua model sehingga perbandingan dapat dilakukan secara berpasangan.

\subsection{Fitur Teknikal sebagai Transformasi Riwayat}

Indikator teknikal meringkas riwayat harga menjadi representasi yang dapat dipelajari oleh model. \emph{Relative Strength Index} (RSI) menggambarkan keseimbangan kenaikan dan penurunan harga, sedangkan \emph{Moving Average Convergence Divergence} (MACD) menggambarkan perbedaan tren eksponensial jangka pendek dan jangka panjang \parencite{wilder1978new,murphy1999technical}. Keduanya merupakan transformasi deterministik dari riwayat harga, sehingga manfaatnya dinilai dari kemampuan representasi tersebut dalam membantu klasifikasi.

Untuk periode RSI $p$, perubahan harga dinyatakan sebagai $d_t=C_t-C_{t-1}$, kenaikan sebagai $g_t=\max(d_t,0)$, dan penurunan sebagai $\ell_t=\max(-d_t,0)$. Setelah rata-rata awal dari $p$ perubahan, penghalusan Wilder dihitung sebagai berikut.
\begin{align}
\overline g_t&=\frac{(p-1)\overline g_{t-1}+g_t}{p}, &
\overline \ell_t&=\frac{(p-1)\overline \ell_{t-1}+\ell_t}{p},\\
RS_t&=\frac{\overline g_t}{\overline \ell_t}, &
RSI_t&=100-\frac{100}{1+RS_t}.
\end{align}
Simbol $p$ menyatakan panjang periode, sedangkan $\overline g_t$ dan $\overline \ell_t$ masing-masing menyatakan rata-rata kenaikan dan penurunan yang dihaluskan pada bar ke-$t$. Jika kedua rata-rata bernilai nol, RSI ditetapkan 50. Jika hanya $\overline \ell_t$ yang nol, RSI ditetapkan 100; jika hanya $\overline g_t$ yang nol, RSI ditetapkan 0.

Rata-rata bergerak eksponensial atau \emph{exponential moving average} (EMA) dengan periode $p$ menggunakan faktor penghalusan $\alpha_p=2/(p+1)$. Perhitungan EMA, MACD, garis sinyal, dan histogram ditunjukkan sebagai berikut.
\begin{align}
\operatorname{EMA}_{p,t}&=\alpha_p C_t+(1-\alpha_p)\operatorname{EMA}_{p,t-1},\\
\operatorname{MACD}_t&=\operatorname{EMA}_{12,t}-\operatorname{EMA}_{26,t},\\
\operatorname{Signal}_t&=\operatorname{EMA}_{9}(\operatorname{MACD})_t,\\
\operatorname{Histogram}_t&=\operatorname{MACD}_t-\operatorname{Signal}_t.
\end{align}
Pada persamaan tersebut, $\operatorname{EMA}_{p,t}$ adalah EMA periode $p$ pada bar ke-$t$, $\operatorname{Signal}_t$ adalah EMA sembilan periode dari MACD, dan $\operatorname{Histogram}_t$ adalah selisih MACD dengan garis sinyal. Nilai awal dan masa pemanasan memengaruhi bagian awal deret, sehingga inisialisasi, normalisasi, dan masa pemanasan akan ditetapkan pada data pelatihan tanpa menggunakan data uji.

\subsection{Fitur Makroekonomi dan Ketertiban Waktu}

Indikator makroekonomi menggambarkan kondisi ekonomi agregat. \emph{Consumer Price Index} (CPI) mengukur perubahan harga sekeranjang barang dan jasa konsumsi, sedangkan CPI inti mengecualikan pangan dan energi. Perubahan bulanan membandingkan indeks dengan bulan sebelumnya, sementara perubahan tahunan membandingkannya dengan bulan yang sama setahun sebelumnya.\footnote{Penjelasan CPI oleh BLS: \url{https://www.bls.gov/cpi/questions-and-answers.htm}.} \emph{Nonfarm Payrolls} (NFP) merujuk pada pekerjaan penggajian nonpertanian; rilis perubahan bulanannya menggambarkan perubahan jumlah pekerjaan. Tingkat pengangguran menyatakan persentase angkatan kerja yang menganggur dan berasal dari survei rumah tangga, sedangkan NFP berasal dari survei perusahaan.\footnote{Catatan teknis statistik ketenagakerjaan BLS: \url{https://www.bls.gov/news.release/empsit.tn.htm}.}

\emph{Federal Open Market Committee} (FOMC) menetapkan arah kebijakan moneter Amerika Serikat, termasuk sasaran suku bunga federal funds.\footnote{Penjelasan FOMC oleh Federal Reserve: \url{https://www.federalreserve.gov/monetarypolicy/fomc.htm}.} Keputusan suku bunga dan publikasi FOMC diperlakukan sebagai peristiwa kalender yang berbeda. Informasi inflasi, ketenagakerjaan, dan kebijakan moneter memberikan konteks bagi pengujian fitur eksternal pada emas. Arah serta besar kontribusi prediktifnya akan diukur melalui eksperimen, tanpa menetapkan bahwa kenaikan suatu indikator selalu diikuti arah harga emas tertentu.

Fitur makroekonomi berasal dari kalender ekonomi terstruktur. Salah satu bentuk utamanya adalah kejutan makro (\emph{macro surprise}), yaitu selisih antara nilai aktual dan prakiraan ($actual - forecast$). Nilai ini menggambarkan perbedaan hasil pengumuman terhadap nilai prakiraan yang tercatat pada kalender. Untuk menjaga urutan waktu kausal dan mencegah kebocoran informasi masa depan, waktu rilis yang telah dikonversi ke UTC harus lebih kecil secara ketat daripada waktu keputusan.

Dalam penelitian ini, konfigurasi K3 dievaluasi menggunakan arsip kalender MetaTrader 5 yang telah dibekukan. Kolom \texttt{forecast} diperlakukan secara operasional sebagai nilai prakiraan yang diasosiasikan vendor dengan peristiwa rilis, sedangkan kolom \texttt{actual} diperlakukan sebagai nilai aktual pada arsip tersebut. Peran semantik kedua kolom ini dibedakan dari verifikasi \emph{point-in-time} independen karena riwayat perubahan nilai tidak tersedia. Oleh sebab itu, kontribusi prediktif K3 ditafsirkan sebagai efektivitas fitur dari snapshot kalender yang digunakan, bukan sebagai pembuktian dampak kausal makroekonomi atau simulasi waktu nyata yang sepenuhnya bebas dari potensi revisi data.

Keadaan tanpa rilis, nilai yang hilang, dan informasi yang kedaluwarsa dibedakan dari kejutan bernilai nol. Kejutan nol berarti nilai aktual sama dengan prakiraan, sedangkan nilai hilang menunjukkan bahwa besarnya kejutan tidak diketahui. Umur rilis dan indikator ketiadaan digunakan untuk mempertahankan perbedaan tersebut. Nilai lama juga tidak diteruskan ketika terdapat rilis yang lebih baru pada keluarga peristiwa yang sama, sehingga informasi lama tidak dipertahankan seolah-olah masih merupakan rilis terkini.

\subsection{Konfigurasi dan Makna Kontribusi Bersyarat}

Tiga konfigurasi fitur dibentuk secara kumulatif. K1 memuat fitur harga, K2 menambahkan fitur teknikal, dan K3 menambahkan fitur makroekonomi dari kalender MetaTrader 5. Untuk metrik $S$, perubahan kinerja akibat penambahan kelompok teknikal dan makroekonomi dihitung dengan Persamaan \ref{eq:delta_teknikal} dan \ref{eq:delta_makro}.
\begin{equation}
\Delta S_{\mathrm{teknikal}}=S_2-S_1
\label{eq:delta_teknikal}
\end{equation}
\begin{equation}
\Delta S_{\mathrm{makro}}=S_3-S_2
\label{eq:delta_makro}
\end{equation}
Simbol $S_1$, $S_2$, dan $S_3$ menyatakan skor K1, K2, dan K3 pada model, metrik, periode uji, label, serta bar uji yang sama. Nilai perubahan dapat positif, nol, atau negatif.

Urutan K1--K2--K3 membuat kontribusi bersifat kondisional. Selisih $S_3-S_2$, misalnya, mengukur perubahan setelah fitur makroekonomi ditambahkan pada fitur harga dan teknikal. Pada analisis utama, setiap konfigurasi akan disetel secara terpisah sehingga selisih juga mencakup respons terhadap penyetelan ulang. Analisis parameter tetap digunakan untuk memeriksa sensitivitas terhadap pilihan tersebut. Pengukuran dilakukan pada tingkat kelompok fitur dan ditafsirkan sebagai perubahan prediktif dalam ruang lingkup penelitian.

\subsection{Regresi Logistik Multinomial}

Regresi Logistik Multinomial memodelkan probabilitas lebih dari dua kelas dan digunakan sebagai model linear \parencite{hosmer2013applied}. Untuk vektor fitur $\mathbf{x}$, probabilitas kelas dihitung dengan fungsi softmax pada Persamaan \ref{eq:softmax}.
\begin{equation}
P(y=k\mid\mathbf{x})=\frac{\exp\left(\beta_{0k}+\boldsymbol{\beta}_k^{\mathsf{T}}\mathbf{x}\right)}{\sum_{j=1}^{K}\exp\left(\beta_{0j}+\boldsymbol{\beta}_j^{\mathsf{T}}\mathbf{x}\right)}
\label{eq:softmax}
\end{equation}
Simbol $K=3$ menyatakan jumlah kelas, $\beta_{0k}$ menyatakan intersep kelas ke-$k$, dan $\boldsymbol{\beta}_k$ menyatakan vektor koefisien kelas ke-$k$. Prediksi ditentukan dari kelas dengan probabilitas terbesar.

Ketidakseimbangan kelas ditangani melalui bobot kelas pada fungsi objektif. Dengan penalti L2, tujuan pelatihan dituliskan pada Persamaan \ref{eq:logistic_objective}.
\begin{equation}
\min_{\boldsymbol{\beta}}\left[-\sum_{i=1}^{n}w_{y_i}\sum_{k=1}^{K}\mathbf{1}(y_i=k)\log P(y_i=k\mid\mathbf{x}_i)+\lambda\sum_{k=1}^{K}\lVert\boldsymbol{\beta}_k\rVert_2^2\right]
\label{eq:logistic_objective}
\end{equation}
Simbol $n$ menyatakan jumlah data pelatihan, $w_{y_i}$ adalah bobot kelas observasi ke-$i$, $\mathbf{1}(y_i=k)$ adalah fungsi indikator, dan $\lambda$ adalah kekuatan regularisasi. Pembobotan mengurangi dominasi kelas yang lebih sering muncul, sedangkan penalti L2 mengendalikan besar koefisien. Pada implementasi, hiperparameter $C$ merupakan kebalikan kekuatan regularisasi sehingga nilai $C$ yang lebih kecil menghasilkan regularisasi lebih kuat. Parameter standardisasi dipelajari hanya dari data pelatihan pada setiap pemisahan latih-validasi atau latih-uji.

\subsection{XGBoost untuk Klasifikasi Multikelas}

XGBoost membentuk prediktor aditif dari sekumpulan pohon keputusan. Vektor skor untuk observasi ke-$i$ dinyatakan pada Persamaan \ref{eq:xgboost_score}.
\begin{equation}
\widehat{\mathbf{s}}_i=\sum_{m=1}^{M}\mathbf{f}_m(\mathbf{x}_i)
\label{eq:xgboost_score}
\end{equation}
Simbol $\widehat{\mathbf{s}}_i$ menyatakan vektor skor untuk tiga kelas, $M$ adalah jumlah tahap boosting, dan $\mathbf{f}_m$ adalah keluaran kelompok pohon pada tahap ke-$m$, termasuk pengaruh laju belajar. Objektif pelatihan menggabungkan kerugian klasifikasi dan penalti kompleksitas pohon, seperti jumlah daun dan besar bobot daun \parencite{chen2016xgboost}.\footnote{Dokumentasi model XGBoost: \url{https://xgboost.readthedocs.io/en/stable/tutorials/model.html}.}

Pada klasifikasi multikelas, vektor skor diubah menjadi probabilitas dan kelas dengan probabilitas terbesar dipilih sebagai prediksi. Struktur pohon dapat menangkap ambang nonlinear dan interaksi antarfitur. Fleksibilitas ini dikendalikan melalui hiperparameter kedalaman, laju belajar, jumlah tahap boosting, subsampling, dan regularisasi yang dipilih pada validasi internal. Dalam penelitian ini, XGBoost digunakan sebagai model nonlinear berbasis pohon.

\subsection{Evaluasi Kronologis dan Pencegahan Kebocoran}

Evaluasi deret waktu mengikuti urutan penggunaan data: periode yang lebih awal digunakan untuk pelatihan dan periode berikutnya digunakan untuk pengujian. Rancangan \emph{nested walk-forward} menggunakan blok validasi internal pada data pelatihan untuk memilih hiperparameter, kemudian model dilatih ulang dan dinilai satu kali pada periode uji utama yang lebih baru. Imputasi, seleksi kolom, penskalaan, bobot kelas, dan penyetelan dipelajari kembali hanya dari bagian pelatihan.

\emph{Purge} menghapus observasi pelatihan yang horizon labelnya bertumpang tindih dengan awal validasi atau pengujian. Sementara itu, embargo merupakan jeda konservatif tambahan setelah batas pemisahan. Riwayat harga dalam jendela fitur telah tersedia pada waktu keputusan, sehingga panjang \emph{purge} ditentukan oleh horizon label. Aturan ini mencegah label pelatihan yang berakhir pada atau setelah keputusan pertama pada blok berikutnya.

Perbandingan konfigurasi dilakukan pada bar, label, periode uji, dan metrik yang sama. Prediktor kelas mayoritas dari data pelatihan digunakan sebagai diagnostik kewajaran. Skor akan dilaporkan per periode uji dan dibandingkan secara berpasangan. Karena periode uji dengan jendela mengembang saling bergantung, ringkasan antarperiode uji ditafsirkan secara deskriptif.

\subsection{Macro-F1 dan Balanced Accuracy}

Matriks kebingungan (\emph{confusion matrix}) merangkum kesesuaian kelas sebenarnya dengan kelas prediksi. Pada Tabel \ref{tab:confusion-tiga-kelas}, $n_{ij}$ adalah jumlah observasi berkelas sebenarnya $i$ yang diprediksi sebagai $j$; $B$, $H$, dan $S$ masing-masing berarti BUY, HOLD, dan SELL. Tabel ini menunjukkan bentuk perhitungan, bukan hasil eksperimen.

\begin{table}[htbp]
\centering\caption{Bentuk matriks kebingungan tiga kelas}\label{tab:confusion-tiga-kelas}
\begin{tabular}{lccc}\hline
Kelas sebenarnya & Prediksi BUY & Prediksi HOLD & Prediksi SELL\\\hline
BUY & $n_{BB}$ & $n_{BH}$ & $n_{BS}$\\
HOLD & $n_{HB}$ & $n_{HH}$ & $n_{HS}$\\
SELL & $n_{SB}$ & $n_{SH}$ & $n_{SS}$\\\hline
\end{tabular}
\end{table}

Untuk setiap kelas $c$, dua kelas lainnya diperlakukan sebagai kelas pembanding. Sebagai contoh, $TP_B=n_{BB}$, $FP_B=n_{HB}+n_{SB}$, dan $FN_B=n_{BH}+n_{BS}$. Dengan demikian, diagonal matriks berisi prediksi yang benar, sedangkan sel di luar diagonal menunjukkan jenis kesalahan klasifikasi.

Untuk kelas $c$, precision, recall, dan F1 dihitung menggunakan Persamaan \ref{eq:precision}--\ref{eq:f1}.
\begin{equation}
\operatorname{Precision}_c=\frac{TP_c}{TP_c+FP_c}
\label{eq:precision}
\end{equation}
\begin{equation}
\operatorname{Recall}_c=\frac{TP_c}{TP_c+FN_c}
\label{eq:recall}
\end{equation}
\begin{equation}
F1_c=\frac{2TP_c}{2TP_c+FP_c+FN_c}
\label{eq:f1}
\end{equation}
Simbol $TP_c$, $FP_c$, dan $FN_c$ masing-masing menyatakan jumlah positif benar, positif salah, dan negatif salah untuk kelas $c$ \parencite{sokolova2009systematic}. Macro-F1 memberi bobot yang sama kepada ketiga kelas seperti ditunjukkan pada Persamaan \ref{eq:macrof1}.
\begin{equation}
\operatorname{Macro\text{-}F1}=\frac{F1_{\mathrm{BUY}}+F1_{\mathrm{HOLD}}+F1_{\mathrm{SELL}}}{3}
\label{eq:macrof1}
\end{equation}
Jika penyebut F1 suatu kelas bernilai nol, F1 kelas tersebut ditetapkan nol dan penyebut macro-F1 tetap tiga kelas. Konvensi ini memungkinkan kinerja setiap kelas tetap terwakili ketika suatu kelas tidak pernah diprediksi.

Balanced accuracy merupakan rata-rata recall kelas dan dihitung dengan Persamaan \ref{eq:balancedacc}.
\begin{equation}
\begin{split}
\operatorname{Balanced\ Accuracy} &= \frac{1}{3} \big( \operatorname{Recall}_{\mathrm{BUY}} \\
&\quad + \operatorname{Recall}_{\mathrm{HOLD}} + \operatorname{Recall}_{\mathrm{SELL}} \big).
\end{split}
\label{eq:balancedacc}
\end{equation}
Jika suatu kelas aktual tidak muncul pada periode uji, recall kelas tersebut dan balanced accuracy tiga kelas tidak terdefinisi. Nilainya tidak diganti dengan nol. Macro-F1 tetap dihitung menggunakan konvensi penyebut tiga kelas, sehingga kedua metrik digunakan secara terpisah untuk membaca perubahan kinerja.
